
\chapter{Logic Gates}
\label{chap:logic-gates}

\chapterintro{Logic gates are the smallest conceptual building blocks of digital computers. This chapter shows how simple true-or-false operations can be combined into circuits that perform computation.}

\learningobjectives{\begin{itemize}\item Explain Boolean values.\item Describe NOT, AND, OR, XOR, NAND, and NOR gates.\item Read basic truth tables.\item Understand how gates lead to arithmetic circuits.\end{itemize}}

\section{Boolean Values}
A Boolean value has only two possible states: true or false. In computing, these are usually represented as \(1\) and \(0\).
\begin{definition}A Boolean value is a value that can be either true or false.\end{definition}

\section{Truth Tables}
A truth table lists every possible input and the output produced by a logical operation.

\section{The NOT Gate}
The NOT gate reverses a Boolean value.
\begin{center}\begin{tabular}{cc}\toprule Input & Output \\\midrule 0 & 1\\ 1 & 0\\\bottomrule\end{tabular}\end{center}

\section{The AND Gate}
The AND gate outputs \(1\) only when both inputs are \(1\).
\begin{center}\begin{tabular}{ccc}\toprule A & B & A AND B \\\midrule 0&0&0\\0&1&0\\1&0&0\\1&1&1\\\bottomrule\end{tabular}\end{center}

\section{The OR Gate}
The OR gate outputs \(1\) when at least one input is \(1\).

\section{Python Example}
\begin{lstlisting}
a = True
b = False
print(not a)
print(a and b)
print(a or b)
\end{lstlisting}

\section{Exercise}
\begin{exercise}Create the truth table for XOR, where the output is true only when the two inputs are different.\end{exercise}
\begin{solution}The XOR output is 0 for 00, 1 for 01, 1 for 10, and 0 for 11.\end{solution}

\section{Summary}
Logic gates transform binary inputs into binary outputs. By combining gates, engineers build adders, memory, CPUs, and eventually the hardware that runs AI systems.
